% !TEX root =./ECE448_Practice_main.tex
%
% L16 practice set (frequency hopping). Every question is a step in the L16
% link: the hop synthesizer's settings, hop timing, the FCC dwell rule,
% jamming, and hop-sequence synchronization.
% Blank copy: \begin{solution}[h] reserves h of work space.
% SOLUTIONS copy (\iskey=1) prints the red worked solutions.

\fullwidth{\textbf{Learning Objective 6.5: } I can implement FHSS transmission and
reception, including hop-sequence synchronization.
\\[4pt]\normalfont\small Show your work. Write a \textbf{Documentation} line at the
end (\textbf{None} if you did not collaborate).}

\begin{questions}

%%%%%%%%%%%%%%%%%%%%%%%%%%%%%%%%%%%%%%%%%
\question \textbf{The hop synthesizer.} A hopper like the L16 transmitter runs at
$f_s = 128$ kHz. A Vector Source lists the hop frequencies, a Repeat block holds each
one for one hop, and a VCO (complex) with sensitivity $2\pi$ turns each value, in
hertz, into a tone. The hopper uses 6 channels spaced 12 kHz apart and centered on
0 Hz, and each hop lasts 125 ms.
\begin{parts}

	\begin{minipage}[t]{0.58\linewidth}
		\part Find the frequency of each channel, $k = 0$ to 5.
		\begin{solution}[1.0in]
		Centered on 0 Hz, channel $k$ sits at
		\eq{ f_k = 12\ \text{kHz} \times (k - 2.5) }
		which gives $-30$, $-18$, $-6$, $+6$, $+18$, and $+30$ kHz.
		\end{solution}
	\end{minipage}\hfill%
	\begin{minipage}[t]{0.37\linewidth}
		\raggedleft \vspace{12pt}
		\ansbox[1.4in]{$\pm 6$, $\pm 18$, $\pm 30$ kHz}
	\end{minipage}

	\begin{minipage}[t]{0.58\linewidth}
		\part Find the Repeat block's interpolation, and the time for one cycle that
		visits every channel once.
		\begin{solution}[1.0in]
		\eq{ 128{,}000 \times 0.125\ \text{s} &= 16{,}000\ \text{samples} \\
		6 \times 125\ \text{ms} &= 750\ \text{ms} }
		\end{solution}
	\end{minipage}\hfill%
	\begin{minipage}[t]{0.37\linewidth}
		\raggedleft \vspace{12pt}
		Repeat: \ansbox[1.0in]{16{,}000}

		\vspace{4pt}
		cycle: \ansbox[1.0in]{750 ms}
	\end{minipage}

	\part The FSK signal on each hop is about 4 kHz wide. In words, explain why a
	channel spacing of 3 kHz would not work, and what the receiver's low-pass filter
	after de-hopping does with the other channels at 12 kHz.
		\begin{solution}[1.1in]
		At 3 kHz spacing each 4 kHz-wide signal would overlap its neighbors, so a
		receiver tuned to one channel would also pick up energy from the next one, and a
		jammer on one channel would spill into two. At 12 kHz the channels do not overlap:
		after de-hopping, the wanted hop sits at 0 Hz and every other channel sits 12 kHz
		or more away, outside a low-pass filter of about $\pm 2$ kHz, which removes them.
		\end{solution}

\end{parts}

\newpage
%%%%%%%%%%%%%%%%%%%%%%%%%%%%%%%%%%%%%%%%%
\question \textbf{Hop timing.} A hopper sends one packet per hop, as the L16 link does:
2 guard bytes, the 12-byte access code and header, the payload, and 2 more
guard bytes. It runs at $R_s = 2000$ baud, with 16 channels.
\begin{parts}

	\begin{minipage}[t]{0.58\linewidth}
		\part For a 26-byte payload, find the bits per hop, the time on one hop, and the
		hop rate.
		\begin{solution}[1.2in]
		\eq{ n &= 8 \times (2 + 12 + 26 + 2) = 336\ \text{bits} \\
		T_\text{hop} &= \dfrac{336}{2000\ \text{baud}} = 168\ \text{ms} \\
		\text{hop rate} &= \dfrac{1}{0.168\ \text{s}} = 5.95\ \text{hops/s} }
		\end{solution}
	\end{minipage}\hfill%
	\begin{minipage}[t]{0.37\linewidth}
		\raggedleft \vspace{12pt}
		bits: \ansbox[1.0in]{336}

		\vspace{4pt}
		hop: \ansbox[1.0in]{168 ms}

		\vspace{4pt}
		rate: \ansbox[1.0in]{5.95 hops/s}
	\end{minipage}

	\begin{minipage}[t]{0.58\linewidth}
		\part Find the time for one cycle of the 16 channels, and the share of each hop
		that carries the payload.
		\begin{solution}[1.0in]
		\eq{ T_\text{cycle} &= 16 \times 168\ \text{ms} = 2.69\ \text{s} \\
		\dfrac{8 \times 26}{336} &= 0.62 = 62\% }
		\end{solution}
	\end{minipage}\hfill%
	\begin{minipage}[t]{0.37\linewidth}
		\raggedleft \vspace{12pt}
		cycle: \ansbox[1.0in]{2.69 s}

		\vspace{4pt}
		payload: \ansbox[1.0in]{62\%}
	\end{minipage}

	\begin{minipage}[t]{0.58\linewidth}
		\part How long is each guard, in milliseconds?
		\begin{solution}[0.7in]
		Each guard is 16 bits, and at 2000 baud it lasts $16/2000 = 8$ ms.
		\end{solution}
	\end{minipage}\hfill%
	\begin{minipage}[t]{0.37\linewidth}
		\raggedleft \vspace{12pt}
		\ansbox[1.1in]{8 ms}
	\end{minipage}

	\part In words, explain why the hopper makes each packet exactly one hop long
	instead of letting packets run across hop boundaries.
		\begin{solution}[1.0in]
		If a packet ran across a boundary, part of it would be on one channel and part on
		the next, so a jammer on either channel, or a receiver whose timing is a little
		off, would lose the whole packet. With one packet per hop the boundaries fall in
		the guard bits between packets, a lost hop loses exactly one packet, and the
		receiver can count packets to check it is on the right hops.
		\end{solution}

\end{parts}

\newpage
%%%%%%%%%%%%%%%%%%%%%%%%%%%%%%%%%%%%%%%%%
\question \textbf{The FCC dwell rule.} Among its limits on power and bandwidth, FCC Part 15.247 lets a narrowband hopper at
915 MHz transmit without a license only if it hops over at least 50 channels, uses every
channel equally, and stays on any one channel no more than 0.4 s in any 20 s. The L16
link hops over 8 channels with 240 ms hops, each channel once per cycle.
\begin{parts}

	\begin{minipage}[t]{0.58\linewidth}
		\part Find how long the L16 link spends on one channel in 20 s.
		\begin{solution}[1.2in]
		The link makes $20 / 0.240 = 83.3$ hops in 20 s, and each channel gets one in 8:
		\eq{ \dfrac{83.3}{8} \times 0.240\ \text{s} = 2.5\ \text{s} }
		which is more than six times the 0.4 s limit.
		\end{solution}
	\end{minipage}\hfill%
	\begin{minipage}[t]{0.37\linewidth}
		\raggedleft \vspace{12pt}
		\ansbox[1.1in]{2.5 s}
	\end{minipage}

	\begin{minipage}[t]{0.58\linewidth}
		\part Set \texttt{nch} to 50, with the same 240 ms hops and each channel once per
		cycle. Find how long the link now spends on one channel in 20 s.
		\begin{solution}[1.0in]
		\eq{ \dfrac{83.3}{50} \times 0.240\ \text{s} = 0.4\ \text{s} }
		which just meets the limit.
		\end{solution}
	\end{minipage}\hfill%
	\begin{minipage}[t]{0.37\linewidth}
		\raggedleft \vspace{12pt}
		\ansbox[1.1in]{0.4 s}
	\end{minipage}

	\part The rule also limits each single visit: a hop longer than 0.4 s breaks it
	even with 50 channels. In words, explain why a permutation sequence, like the
	L16 link's, makes the equal-use rule easy to meet, and what a sequence drawn
	independently at random each hop would do instead.
		\begin{solution}[1.1in]
		A permutation visits every channel exactly once per cycle, so every channel gets
		exactly the same time in every cycle, and the average use is equal by
		construction. A channel drawn at random each hop is equal only on average: over a
		short window some channels come up more often than others, and one could exceed
		0.4 s in a 20 s window even when the long-run average is fine. It could also
		repeat a channel on consecutive hops.
		\end{solution}

\end{parts}

\newpage
%%%%%%%%%%%%%%%%%%%%%%%%%%%%%%%%%%%%%%%%%
\question \textbf{Jamming.} The L16 link hops over 8 channels, channel $k$ at
$10\ \text{kHz} \times (k - 3.5)$, in the order 7, 2, 0, 6, 4, 1, 3, 5. A jammer tone sits
at $+5$ kHz, on channel 4. On each hop, Hop Sync multiplies the band by a tone at minus
that hop's channel frequency, and the low-pass filter after it keeps $\pm 2$ kHz.
\begin{parts}

	\begin{minipage}[t]{0.58\linewidth}
		\part Find where the jammer appears after Hop Sync on the hop to channel 7, and on
		the hop to channel 4.
		\begin{solution}[1.0in]
		Hop Sync subtracts the hop's frequency. Channel 7 is at $+35$ kHz, so the jammer
		lands at $5 - 35 = -30$ kHz; channel 4 is at $+5$ kHz, so it lands at 0 Hz.
		\end{solution}
	\end{minipage}\hfill%
	\begin{minipage}[t]{0.37\linewidth}
		\raggedleft \vspace{12pt}
		ch 7: \ansbox[1.0in]{$-30$ kHz}

		\vspace{4pt}
		ch 4: \ansbox[1.0in]{0 Hz}
	\end{minipage}

	\begin{minipage}[t]{0.58\linewidth}
		\part On which hops of the cycle does the jammer get through the low-pass filter?
		\begin{solution}[1.0in]
		Only on the hop to channel 4. On every other hop the jammer lands at a multiple of
		10 kHz away from 0 Hz, at least 10 kHz, far outside $\pm 2$ kHz.
		\end{solution}
	\end{minipage}\hfill%
	\begin{minipage}[t]{0.37\linewidth}
		\raggedleft \vspace{12pt}
		\ansbox[1.1in]{1 hop in 8}
	\end{minipage}

	\part A cadet turns the jammer amplitude up from 1 to 10, and the Packets label stays
	near 88\%. Explain in words why the extra power takes no more packets.
		\begin{solution}[1.0in]
		The jammer already takes every hop on channel 4 at amplitude 1. On the other 7 hops
		of the cycle the filter removes it however strong it is, so raising its power
		cannot reach those hops; only a jammer on more channels could.
		\end{solution}

	\part Bluetooth's adaptive frequency hopping lets the receiver mark channels that
	keep losing packets as bad and drop them from the sequence. A hopper with 20 channels
	has 3 of them jammed. In words, explain what happens to the packets delivered when it
	drops the 3 jammed channels, and why the jammer cannot simply follow.
		\begin{solution}[1.2in]
		With the 3 jammed channels dropped, the hopper uses the other 17, none of which is
		jammed, so every packet arrives again. The jammer does not know the sequence, so it
		cannot tell which channels are now in use except by listening; to hit the hops it
		would have to move to new channels and cover them too, and then the receiver drops
		those. Each channel the hopper drops makes it a little easier to find. Bluetooth's AFH must
		keep at least 20 channels, a floor set by the radio regulations.
		\end{solution}

\end{parts}

\newpage
%%%%%%%%%%%%%%%%%%%%%%%%%%%%%%%%%%%%%%%%%
\question \textbf{Synchronization.} The L16 link hops over 8 channels with key 448, in
the order 7, 2, 0, 6, 4, 1, 3, 5, with 240 ms hops and 16 guard bits, 16 ms, at each end
of every hop. Its receiver parks on the sync channel, channel 7, until hop 0 arrives
there.
\begin{parts}

	\begin{minipage}[t]{0.58\linewidth}
		\part A receiver with the right timing tunes one hop behind the transmitter. List
		the transmitter's and the receiver's channels for hops 1 to 4, and find how many
		packets arrive in a full cycle.
		\begin{solution}[1.2in]
		On hop $k$ the transmitter is on entry $k$ of the list and the receiver on entry
		$k-1$: hop 1 is 2 and 7, hop 2 is 0 and 2, hop 3 is 6 and 0, hop 4 is 4 and 6.
		Consecutive entries of a permutation are never equal, so they never match, and no
		packet arrives in any cycle.
		\end{solution}
	\end{minipage}\hfill%
	\begin{minipage}[t]{0.37\linewidth}
		\raggedleft \vspace{12pt}
		\ansbox[1.1in]{0}
	\end{minipage}

	\begin{minipage}[t]{0.58\linewidth}
		\part A receiver starts and parks on channel 7 at 0.70 s after the start of a
		hop 0. How long does it wait before it locks?
		\begin{solution}[0.9in]
		The next hop 0 starts one cycle after the last, at $8 \times 0.240 = 1.92$ s:
		\eq{ 1.92 - 0.70 = 1.22\ \text{s} }
		\end{solution}
	\end{minipage}\hfill%
	\begin{minipage}[t]{0.37\linewidth}
		\raggedleft \vspace{12pt}
		\ansbox[1.1in]{1.22 s}
	\end{minipage}

	\part After locking, the receiver changes channel 10 ms late on every hop. A
	teammate's receiver changes channel 20 ms early. Explain what each one misses and
	whether its packets arrive.
		\begin{solution}[1.2in]
		The late receiver spends the first 10 ms of each hop on the previous channel, so it
		misses 10 of the 16 leading guard bits; the access code, header, and payload all
		arrive, and so does every packet. The early receiver leaves 20 ms before the hop
		ends, so it misses the 16 trailing guard bits and the last 4 bits of the newline;
		the payload is wrong, and almost no packet arrives intact.
		\end{solution}

	\part In words, explain why parking on one channel tells the receiver which hop the
	transmitter is on, and what would go wrong if the sequence visited channel 7 twice in
	each cycle.
		\begin{solution}[1.1in]
		The sequence is a permutation, so the transmitter is on channel 7 only on hop 0;
		power appearing there means hop 0 has just started, which gives both the hop
		number and the hop boundary. If channel 7 appeared twice, the receiver could not
		tell which of the two visits it heard, and half the time it would lock to the
		wrong hop number, off by several hops, and receive nothing until it parked again.
		\end{solution}

\end{parts}

\vspace{1.5em}
\noindent\textbf{Documentation:} \rule[-2pt]{0.6\linewidth}{0.4pt}

\end{questions}
